\documentclass[a4paper,11pt]{article}
\usepackage[utf8]{inputenc}
\usepackage[french]{babel}
\usepackage[T1]{fontenc}
\usepackage{amsmath,amssymb,amsthm}
\usepackage{geometry}
\geometry{margin=2cm}
\usepackage{enumitem}
\usepackage{hyperref}
\usepackage{booktabs}
\usepackage{array}
\usepackage{xcolor}
\usepackage{etoolbox}
\AtBeginDocument{\shorthandoff{;:!?}}
\usepackage{listings}
\usepackage{titlesec}
\usepackage{fancyhdr}
\usepackage{lastpage}
\usepackage{graphicx}

% ---- Style des titres ----
\titleformat{\section}{\large\bfseries\color{blue!60!black}}{}{0pt}{}[\vspace{-0.5ex}\rule{\textwidth}{0.4pt}]
\titleformat{\subsection}{\bfseries\color{blue!40!black}}{}{0pt}{}
\titleformat{\subsubsection}{\itshape\color{blue!30!black}}{}{0pt}{}

% ---- Style listings ----
\lstdefinestyle{monstyle}{
  frame=single,
  numbers=left,
  numbersep=5pt,
  breaklines=true,
  basicstyle=\small\ttfamily,
  backgroundcolor=\color{gray!5},
  rulecolor=\color{gray!40},
  columns=fullflexible,
  keepspaces=true,
  showstringspaces=false
}
\lstset{style=monstyle}

% ---- En-têtes / pieds ----
\pagestyle{fancy}
\fancyhf{}
\fancyhead[L]{\small STA211 -- Pré-traitement}
\fancyhead[R]{\small Fiche récapitulative}
\fancyfoot[C]{\thepage/\pageref{LastPage}}
\renewcommand{\headrulewidth}{0.4pt}

% ---- Theorem-like environment ----
\theoremstyle{definition}
\newtheorem*{definition}{Définition}
\theoremstyle{remark}
\newtheorem*{remark}{Remarque}
\newtheorem*{important}{À retenir}

\title{\textbf{Fiche récapitulative -- STA211}\\[4pt]
\large Exploitation des méthodes exploratoires\\pour le pré-traitement des données}
\author{}
\date{10 juin 2026}

\begin{document}
\maketitle
\thispagestyle{empty}

\begin{abstract}
Cette fiche synthétise le cours d'Exploitation des méthodes exploratoires pour le
pré-traitement des données (V. Audigier, N. Niang). Le pré-traitement transforme
les données brutes en un format adapté à l'analyse. Les méthodes exploratoires
aident à comprendre la structure des données avant d'appliquer des modèles
prédictifs.
\end{abstract}

\vfill
\tableofcontents
\newpage

% ===================================================================
\section{Présentation}
% ===================================================================

Le pré-traitement est une étape cruciale de tout processus d'analyse de données.
Les données brutes contiennent souvent des imperfections~: valeurs aberrantes,
données manquantes, échelles hétérogènes, variables redondantes, etc. Ignorer
ces problèmes conduit à des résultats biaisés ou peu robustes.

Les méthodes exploratoires jouent un double rôle~:
\begin{enumerate}
  \item \textbf{Diagnostic}~: mettre en évidence les caractéristiques des données
    (distributions, corrélations, atypiques, manquantes).
  \item \textbf{Transformation}~: corriger ou adapter les données pour les rendre
    compatibles avec les méthodes d'analyse ultérieures.
\end{enumerate}

Questions clés abordées dans ce cours~:
\begin{itemize}
  \item Comment transformer les variables (quantitatives et qualitatives)~?
  \item Comment détecter et traiter les points atypiques~?
  \item Comment gérer les données manquantes~?
  \item Comment réduire la dimension des données~?
  \item Comment traiter les données volumineuses~?
\end{itemize}

% ===================================================================
\section{Transformation des variables quantitatives}
% ===================================================================

\subsection{Discrétisation}

La discrétisation transforme une variable continue en classes discrètes.
Utile pour~: aplicabilité de certains algorithmes, interprétabilité,
robustesse au bruit.

\begin{description}[style=nextline]
  \item[Règles métier] Découpage selon des seuils ayant un sens métier.
    Exemple~: âge en classes~: $[0,18[$, $[18,25[$, $[25,45[$, $[45,65[$,
    $65+$. Simple, pertinent mais subjectif.
  \item[Quantiles] Classes d'effectifs égaux (déciles, quartiles).
    Exemple~: âge en 3 quantiles~: $[18,32[$, $[32,51[$, $[51,95]$.
    Objectif mais peut créer des frontières artificielles.
  \item[CAH / $k$-means] On regroupe les valeurs par clustering univarié
    puis on utilise les clusters comme catégories. S'adapte à la structure
    naturelle des données.
  \item[MDLPC] \textit{Minimum Description Length Principle Cut}~:
    discrétisation supervisée qui minimise la perte d'information vis-à-vis
    d'une cible $Y$. Choisit les coupures qui maximisent le gain d'information.
\end{description}

\begin{important}
  La discrétisation fait perdre de l'information (perte de granularité).
  À utiliser avec parcimonie, et valider l'impact sur l'analyse aval.
\end{important}

\subsection{Linéarité}

De nombreux modèles (régression linéaire, ACP) supposent des relations
linéaires entre variables. Une transformation peut linéariser la relation.

\textbf{Détection}~:
\begin{itemize}
  \item Coefficient de Pearson $r$~: mesure la corrélation linéaire.
  \item Coefficient de Spearman $\rho$~: mesure la corrélation de rang
    (monotone, pas forcément linéaire).
  \item Si $|r| \approx |\rho|$~: relation approximativement linéaire.
  \item Si $|\rho| > |r|$~: relation monotone mais non linéaire → besoin
    de transformation.
\end{itemize}

\textbf{Transformation de Box-Cox}~:
\[
y^{(\lambda)} =
\begin{cases}
  \dfrac{y^\lambda - 1}{\lambda} & \text{si } \lambda \neq 0,\\[6pt]
  \log(y)                         & \text{si } \lambda = 0.
\end{cases}
\]
Le paramètre $\lambda$ est estimé par maximum de vraisemblance (ou
\textit{profile likelihood}). On cherche $\lambda$ tel que $y^{(\lambda)}$
soit le plus linéaire possible par rapport aux autres variables.

\subsection{Normalité}

De nombreux tests et modèles supposent la normalité des résidus ou des
variables.

\textbf{Indicateurs de non-normalité}~:
\begin{itemize}
  \item Asymétrie (\textit{skewness})~:
    \[
    \gamma_1 = \mathbb{E}\!\left[\left(\frac{X-\mu}{\sigma}\right)^{\!3}\right].
    \]
    $\gamma_1 = 0$ pour une normale. $\gamma_1 > 0$~: queue à droite.
    $\gamma_1 < 0$~: queue à gauche.
  \item Aplatissement (\textit{kurtosis})~:
    \[
    \gamma_2 = \mathbb{E}\!\left[\left(\frac{X-\mu}{\sigma}\right)^{\!4}\right] - 3.
    \]
    $\gamma_2 = 0$ pour une normale. $\gamma_2 > 0$~: distribution
    leptikurtique (pointue). $\gamma_2 < 0$~: platykurtique (aplatie).
  \item Test de Shapiro-Wilk, test de Jarque-Bera.
  \item QQ-plot~: quantiles empiriques vs quantiles théoriques.
    L'alignement des points sur la diagonale indique la normalité.
\end{itemize}

\textbf{Transformations correctives}~:
\begin{itemize}
  \item Logarithme~: $\log(x)$ pour des données strictement positives
    avec asymétrie droite.
  \item Racine carrée~: $\sqrt{x}$ pour des données de comptage.
  \item Box-Cox (cf.~ci-dessus)~: généralisation.
  \item Yeo-Johnson~: analogue à Box-Cox mais accepte les valeurs
    négatives.
\end{itemize}

% ===================================================================
\section{Transformation des variables qualitatives}
% ===================================================================

\subsection{Regroupement de modalités}

Les variables qualitatives comportent parfois un grand nombre de modalités,
dont certaines très rares. Le regroupement vise à~:
\begin{itemize}
  \item Éviter les indicateurs binaires de faible effectif (instables).
  \item Améliorer la parcimonie et l'interprétabilité.
\end{itemize}

\textbf{Méthodes}~:
\begin{enumerate}
  \item Expertise métier~: regroupement basé sur la connaissance du domaine.
  \item Clustering des modalités~: on représente chaque modalité par son
    profil (e.g.\ via ACM), puis on regroupe les modalités proches.
  \item Fusion supervisée~: pour une cible $Y$ donnée, on regroupe les
    modalités dont la moyenne de $Y$ est similaire (test de Student ou
    ANOVA).
\end{enumerate}

\subsection{Analyse factorielle}

L'analyse factorielle (ACM, AFDM) permet de convertir des variables
qualitatives en représentations numériques continues. Les coordonnées
des individus sur les axes factoriels deviennent de nouvelles variables
quantitatives.

\begin{itemize}
  \item \textbf{ACM (MCA)}~: pour données qualitatives uniquement.
    Transforme le tableau disjonctif complet en composantes continues.
  \item \textbf{AFDM (FAMD)}~: pour données mixtes (quantitatives +
    qualitatives). Combine ACP (pour les quantitatives) et ACM (pour
    les qualitatives).
\end{itemize}

\begin{important}
  L'analyse factorielle n'est pas seulement une méthode de visualisation~:
  c'est aussi un outil de pré-traitement pour remplacer des variables
  brutes par des composantes synthétiques.
\end{important}

% ===================================================================
\section{Détection des points atypiques}
% ===================================================================

\subsection{Méthodes univariées}

\textbf{Z-score}~:
\[
z_i = \frac{x_i - \bar{x}}{s},
\]
où $\bar{x}$ est la moyenne et $s$ l'écart-type. On considère comme
atypique une observation avec $|z_i| > 3$ (seuil de la normale à 99,7~\%).
Limite~: la moyenne et l'écart-type sont sensibles aux valeurs extrêmes
(effet de masquage).

\textbf{Intervalle interquartile (IQR)}~:
\[
\text{IQR} = Q_3 - Q_1.
\]
Sont atypiques les valeurs $< Q_1 - 1{,}5 \times \text{IQR}$ ou
$> Q_3 + 1{,}5 \times \text{IQR}$. Plus robuste que le Z-score car
basé sur les quartiles.

\subsection{Méthodes multivariées}

\textbf{Distance de Mahalanobis}~:
\[
D_i^2 = (\mathbf{x}_i - \boldsymbol{\mu})^\top \boldsymbol{\Sigma}^{-1}
        (\mathbf{x}_i - \boldsymbol{\mu}).
\]
Sous hypothèse de normalité multivariée, $D_i^2 \sim \chi^2_p$ avec $p$
le nombre de variables.
On considère atypiques les individus dont la $p$-valeur associée est
$< 0{,}001$ (ou autre seuil).

En pratique, $\boldsymbol{\mu}$ et $\boldsymbol{\Sigma}$ sont estimés
sur l'échantillon, mais ces estimateurs sont sensibles aux atypiques.
On utilise alors l'estimateur MCD (\textit{Minimum Covariance
Determinant}) qui minimise le déterminant de la matrice de covariance
sur un sous-ensemble robuste.

\textbf{Isolation Forest}~\cite{liu2008isolation} :
\begin{itemize}
  \item Algorithme non supervisé basé sur des arbres de décision aléatoires.
  \item Principe~: on isole chaque observation en partitionnant
    aléatoirement l'espace des variables.
  \item Un point atypique est isolé en moins de coupes (chemin plus court
    dans l'arbre) qu'un point normal.
\end{itemize}
Score d'anomalie~:
\[
s(x, n) = 2^{-\frac{\mathbb{E}[h(x)]}{c(n)}}, \qquad
c(n) = 2H(n-1) - \frac{2(n-1)}{n},
\]
où $h(x)$ est la longueur du chemin dans l'arbre, $c(n)$ un facteur de
normalisation et $H(i)$ le $i$-ième nombre harmonique.
$s \to 1$~: atypique, $s \to 0$~: normal.

% ===================================================================
\section{Gestion des données manquantes}
% ===================================================================

\subsection{Mécanismes de manquance}

\begin{description}[style=nextline]
  \item[MCAR] (\textit{Missing Completely At Random})~: la probabilité
    de manquance ne dépend ni des données observées ni des données
    manquantes. Exemple~: questionnaire perdu aléatoirement.
  \item[MAR] (\textit{Missing At Random})~: la probabilité de manquance
    dépend des données observées mais pas des données manquantes.
    Exemple~: les hommes répondent moins à une question sur le poids.
  \item[MNAR] (\textit{Missing Not At Random})~: la probabilité de
    manquance dépend des données manquantes elles-mêmes.
    Exemple~: les personnes à haut revenu ne déclarent pas leur revenu.
\end{description}

\begin{important}
  Le mécanisme de manquance est une hypothèse non testable
  (car on n'observe pas les données manquantes). Il faut justifier
  son choix par des arguments métier.
\end{important}

\subsection{Approches par suppression}

\begin{itemize}
  \item \textbf{Listwise deletion (analyse des cas complets)}~:
    on supprime toute ligne comportant au moins une valeur manquante.
    Simple mais~: (i) perte de puissance, (ii) biais sauf si MCAR.
  \item \textbf{Pairwise deletion}~: on utilise toutes les observations
    disponibles pour chaque calcul (moyenne, corrélation). La matrice de
    covariance peut ne pas être définie positive.
\end{itemize}

\subsection{Imputation simple}

\begin{itemize}
  \item \textbf{Moyenne / médiane}~: remplace la valeur manquante par la
    moyenne (ou médiane) de la variable. Attenue la variance et les
    corrélations.
  \item \textbf{Imputation par régression}~: on prédit la valeur manquante
    à partir des autres variables. Tendance à renforcer les relations
    existantes.
  \item \textbf{Hot-deck ($k$-NN)}~: on trouve les $k$ plus proches voisins
    complets et on impute par leur moyenne (ou un tirage aléatoire).
\end{itemize}

\subsection{Imputation multiple -- MICE}

\textbf{MICE} (\textit{Multiple Imputation by Chained Equations})~:
méthode itérative qui génère $m$ jeux de données imputés, en tenant
compte de l'incertitude d'imputation.

\textbf{Algorithme}~:
\begin{enumerate}
  \item Imputation initiale provisoire (e.g.\ moyenne).
  \item Pour chaque variable $X_j$ contenant des valeurs manquantes~:
    \begin{enumerate}
      \item Régression de $X_j$ sur toutes les autres variables
        (sur les observations où $X_j$ est observée).
      \item Imputation des valeurs manquantes de $X_j$ par tirage
        dans la distribution prédictive (valeur prédite + bruit).
    \end{enumerate}
  \item Répéter l'étape 2 pendant plusieurs cycles (5--10) pour
    stabiliser les imputations.
  \item Répéter tout le processus $m$ fois (typiquement $m=5$ ou
    $m=20$) → $m$ jeux de données complets.
  \item Analyser chaque jeu séparément et combiner les résultats
    avec les règles de Rubin.
\end{enumerate}

\textbf{Règles de Rubin}~:
\[
\bar{Q} = \frac{1}{m}\sum_{j=1}^m \hat{Q}_j,\qquad
B = \frac{1}{m-1}\sum_{j=1}^m (\hat{Q}_j - \bar{Q})^2,\qquad
T = \bar{U} + \left(1 + \frac{1}{m}\right)B.
\]
où $\hat{Q}_j$ est l'estimation du paramètre d'intérêt sur le
$j$-ème jeu imputé, $U_j$ sa variance estimée,
$\bar{U} = \frac{1}{m}\sum U_j$ la variance intra-imputation,
$B$ la variance inter-imputation, et $T$ la variance totale.

\section{Réduction de dimension}

\subsection{Sélection \textit{vs} extraction}

Deux stratégies complémentaires~:
\begin{itemize}
  \item \textbf{Sélection de variables}~: choisir un sous-ensemble des
    variables originales. Avantage~: interprétabilité préservée.
    Inconvénient~: information redondante partiellement perdue.
  \item \textbf{Extraction de caractéristiques}~: créer de nouvelles
    variables synthétiques (composantes, facteurs). Avantage~: capte
    l'information de toutes les variables. Inconvénient~: perte
    d'interprétabilité directe.
\end{itemize}

\subsection{Méthodes d'extraction}

\begin{description}[style=nextline]
  \item[ACP (PCA)] Trouve les directions de variance maximale dans
    l'espace des variables. Les composantes principales sont
    orthogonales. Adaptée aux variables numériques.
  \item[ACM (MCA)] Version de l'ACP pour variables qualitatives.
    Appliquée au tableau disjonctif complet.
  \item[AFDM (FAMD)] Mélange ACP + ACM pour données mixtes.
  \item[PLS (\textit{Partial Least Squares})] Trouve des composantes
    qui maximisent la covariance entre $X$ et $Y$. Supervisée.
\end{description}

\subsection{Méthodes de sélection}

\begin{itemize}
  \item \textbf{Best subset}~: testing all $2^p$ subsets. Impossible
    pour $p$ grand.
  \item \textbf{Sélection pas-à-pas}~: forward (ajout), backward
    (retrait), stepwise (les deux).
  \item \textbf{Régularisation (Lasso, Ridge, Elastic Net)}~:
    \[
    \hat{\boldsymbol{\beta}} = \arg\min_{\boldsymbol{\beta}}
    \| \mathbf{y} - \mathbf{X}\boldsymbol{\beta} \|_2^2
    + \lambda \, \|\boldsymbol{\beta}\|_q^q,
    \]
    avec $q=1$ (Lasso, fait de la sélection), $q=2$ (Ridge, réduit
    sans sélectionner), ou un mélange (Elastic Net).
\end{itemize}

\subsection{Clustering de variables}

Objectif~: identifier des groupes de variables corrélées pour détecter
la redondance.

Démarche~:
\begin{enumerate}
  \item Calculer la matrice de corrélation $\mathbf{R}$.
  \item Définir une distance entre variables, par exemple
    $d(X_j, X_k) = 1 - |r_{jk}|$.
  \item Appliquer une CAH sur les variables (et non sur les individus).
\end{enumerate}

Chaque cluster peut ensuite être résumé par une variable synthétique
(la première composante de l'ACP intra-cluster).

% ===================================================================
\section{Pré-traitement des données volumineuses}
% ===================================================================

Avec des jeux de données massifs, les méthodes classiques deviennent
coûteuses, voire impossibles à exécuter en mémoire.

\textbf{Stratégies}~:
\begin{itemize}
  \item \textbf{Échantillonnage}~: réduire la taille des données.
    \begin{itemize}
      \item Aléatoire simple~: tirage uniforme sans remise.
      \item Systématique~: on sélectionne 1 individu tous les $k$.
      \item Stratifié~: tirage par strates pour préserver la structure.
    \end{itemize}
  \item \textbf{Sous-échantillonnage pour Isolation Forest}~: chaque
    arbre est construit sur un petit sous-échantillon (e.g.\ 256
    individus). Cela améliore la robustesse et réduit le coût.
  \item \textbf{Parallélisation}~: diviser les données en blocs,
    traiter chaque bloc indépendamment, puis combiner les résultats.
  \item \textbf{Algorithmes en ligne (\textit{streaming})}~: mise à
    jour incrémentale des paramètres (e.g.\ SGD, ACP incrémentale).
\end{itemize}

% ===================================================================
\section{Code Python -- Exemples}
% ===================================================================

\begin{lstlisting}[language=Python, caption=Box-Cox et détection d'atypiques]
import numpy as np
from scipy.stats import boxcox
from scipy.spatial.distance import mahalanobis

# --- Transformation Box-Cox ---
data = np.array([1, 2, 3, 5, 10, 20, 50])
transformed, lambda_opt = boxcox(data)
print(f"lambda optimal: {lambda_opt:.3f}")

# --- Distance de Mahalanobis robuste ---
from sklearn.covariance import MinCovDet
X = np.random.randn(100, 5)
mcd = MinCovDet().fit(X)
dist = mcd.mahalanobis(X)
seuil = np.sqrt(mcd.dof_)  # approx chi2(p)
outliers = dist > seuil
print(f"Atteinte robuste Mahalanobis: {outliers.sum()} atypiques")
\end{lstlisting}

\begin{lstlisting}[language=Python, caption=Isolation Forest]
from sklearn.ensemble import IsolationForest
rng = np.random.RandomState(42)
X = 0.3 * rng.randn(100, 2)
X_outliers = rng.uniform(low=-4, high=4, size=(20, 2))
X = np.vstack([X, X_outliers])

model = IsolationForest(contamination=0.1, random_state=rng)
y_pred = model.fit_predict(X)
# -1 = outlier, 1 = normal
print(f"Atteinte Isolation Forest: {(y_pred == -1).sum()} atypiques")
\end{lstlisting}

\begin{lstlisting}[language=Python, caption=Imputation multiple avec IterativeImputer]
from sklearn.experimental import enable_iterative_imputer
from sklearn.impute import IterativeImputer
import pandas as pd

df = pd.DataFrame({
    'X1': [1, 2, np.nan, 4, 5],
    'X2': [5, np.nan, 3, 2, 1],
    'X3': [1, 1, 0, 0, 1]
})
imp = IterativeImputer(max_iter=10, random_state=0)
df_imputed = pd.DataFrame(imp.fit_transform(df),
                          columns=df.columns)
print(df_imputed)
\end{lstlisting}

\begin{lstlisting}[language=Python, caption=ACP et Lasso]
from sklearn.decomposition import PCA
from sklearn.linear_model import LassoCV
from sklearn.preprocessing import StandardScaler

# ACP
scaler = StandardScaler()
X_scaled = scaler.fit_transform(X)
pca = PCA(n_components=2)
X_pca = pca.fit_transform(X_scaled)
print(f"Variance expliquee: {pca.explained_variance_ratio_}")

# Lasso
lasso = LassoCV(cv=5, random_state=0).fit(X_scaled, y)
print(f"Variables selectionnees: {np.where(lasso.coef_ != 0)[0]}")
\end{lstlisting}

% ===================================================================
\section{Code R -- Exemples}
% ===================================================================

\begin{lstlisting}[language=R, caption=Box-Cox et Mahalanobis]
library(MASS)
library(robustbase)

# --- Box-Cox ---
boxcox(lm(mpg ~ ., data = mtcars))

# --- Distance de Mahalanobis robuste ---
data(iris)
X <- iris[, 1:4]
mcd <- covMcd(X)
dist <- mahalanobis(X, center = mcd$center, cov = mcd$cov)
seuil <- qchisq(0.999, df = ncol(X))
outliers <- dist > seuil
cat("Atteinte Mahalanobis robuste:", sum(outliers), "atypiques\n")
\end{lstlisting}

\begin{lstlisting}[language=R, caption=Isolation Forest avec isotree]
library(isotree)
set.seed(42)
iso <- isolation.forest(X, ntrees = 100, sample_size = 256,
                        ndim = ncol(X))
scores <- predict(iso, X)
cat("Score d'anomalie (moy.):", mean(scores), "\n")
\end{lstlisting}

\begin{lstlisting}[language=R, caption=Imputation multiple avec mice]
library(mice)
data(nhanes)
imp <- mice(nhanes, m = 5, method = 'pmm', seed = 123)
# Combinaison des resultats (regression)
fit <- with(imp, lm(bmi ~ age + chl))
pooled <- pool(fit)
summary(pooled)
\end{lstlisting}

\begin{lstlisting}[language=R, caption={ACP, FAMD et Lasso}]
library(FactoMineR)
library(glmnet)

# ACP
res.pca <- PCA(iris[, 1:4], scale.unit = TRUE, graph = FALSE)
print(res.pca$eig[1:2, ])

# FAMD (donnees mixtes)
data(geomorphology)
res.famd <- FAMD(geomorphology, graph = FALSE)

# Lasso
X <- as.matrix(iris[, 1:4])
y <- iris[, 5]
cv.lasso <- cv.glmnet(X, as.numeric(y), alpha = 1, nfolds = 5)
coef(cv.lasso, s = "lambda.1se")
\end{lstlisting}

% ===================================================================
\section*{Questions clés (QCM)}
% ===================================================================

\addcontentsline{toc}{section}{Questions clés (QCM)}

Voici 10 questions pour vérifier votre compréhension. Les réponses
sont fournies ci-dessous.

\begin{enumerate}

  \item Quelle transformation est recommandée pour linéariser une
    relation monotone non linéaire~?
    \begin{enumerate}
      \item Transformation logarithmique
      \item Transformation de Box-Cox
      \item Centrage-réduction
      \item Discrétisation par quantiles
    \end{enumerate}
    \textbf{Réponse~: b}

  \item Le coefficient de Spearman mesure~:
    \begin{enumerate}
      \item La corrélation linéaire entre deux variables
      \item La corrélation de rang (monotone)
      \item La proportion de variance expliquée
      \item La distance entre deux distributions
    \end{enumerate}
    \textbf{Réponse~: b}

  \item Dans la distance de Mahalanobis, la matrice utilisée est~:
    \begin{enumerate}
      \item La matrice des corrélations
      \item L'inverse de la matrice de variance-covariance
      \item La matrice des distances euclidiennes
      \item La matrice identité
    \end{enumerate}
    \textbf{Réponse~: b}

  \item Un point atypique dans Isolation Forest a~:
    \begin{enumerate}
      \item Un chemin long dans l'arbre
      \item Un chemin court dans l'arbre
      \item Un score proche de 0
      \item Une forte contribution à l'inertie
    \end{enumerate}
    \textbf{Réponse~: b}

  \item Le mécanisme MAR signifie que la probabilité de manquance~:
    \begin{enumerate}
      \item Est indépendante de toutes les données
      \item Dépend des données observées uniquement
      \item Dépend des données manquantes uniquement
      \item Dépend des données observées et manquantes
    \end{enumerate}
    \textbf{Réponse~: b}

  \item Dans MICE, le nombre $m$ de jeux de données imputés est
    typiquement~:
    \begin{enumerate}
      \item $m=1$
      \item $m=5$ à $20$
      \item $m=100$ à $500$
      \item $m$ est choisi par validation croisée
    \end{enumerate}
    \textbf{Réponse~: b}

  \item Le Lasso effectue de la sélection de variables car~:
    \begin{enumerate}
      \item Il utilise la norme $\ell_2$
      \item Il utilise la norme $\ell_1$
      \item Il estime les coefficients par MCCO
      \item Il orthogonalise les prédicteurs
    \end{enumerate}
    \textbf{Réponse~: b}

  \item L'AFDM est adaptée aux données~:
    \begin{enumerate}
      \item Uniquement qualitatives
      \item Uniquement quantitatives
      \item Mixtes (qualitatives et quantitatives)
      \item Textuelles
    \end{enumerate}
    \textbf{Réponse~: c}

  \item Le clustering de variables permet de~:
    \begin{enumerate}
      \item Regrouper des individus similaires
      \item Identifier des groupes de variables corrélées
      \item Réduire le nombre d'individus
      \item Détecter des points atypiques
    \end{enumerate}
    \textbf{Réponse~: b}

  \item Pour des données volumineuses, une stratégie courante
    dans Isolation Forest est~:
    \begin{enumerate}
      \item Utiliser tous les individus pour chaque arbre
      \item Sous-échantillonner les individus pour chaque arbre
      \item Augmenter le nombre de variables
      \item Utiliser une ACP en amont
    \end{enumerate}
    \textbf{Réponse~: b}

\end{enumerate}

% ===================================================================
% Bibliographie
% ===================================================================
\begin{thebibliography}{9}

\bibitem{liu2008isolation}
F.T.~Liu, K.M.~Ting, Z.-H.~Zhou.
\newblock Isolation Forest.
\newblock \emph{ICDM}, 2008.

\bibitem{audigier2023}
V.~Audigier, N.~Niang.
\newblock Exploitation des méthodes exploratoires pour le pré-traitement
  des données.
\newblock Cours STA211, 2023.

\bibitem{vanbuuren2018}
S.~van~Buuren.
\newblock \emph{Flexible Imputation of Missing Data}.
\newblock 2nd ed., CRC Press, 2018.

\bibitem{tibshirani1996}
R.~Tibshirani.
\newblock Regression Shrinkage and Selection via the Lasso.
\newblock \emph{JRSS B}, 58(1):267--288, 1996.

\end{thebibliography}

\end{document}
